\chapter{System Architecture}

\section{Overview}
Project Vanish Lite is partitioned into two primary functional halves: the low-level Local Execution System handling daemon creation and restrictions, and the higher-level Web Admin Panel orchestrating monitoring and visualization. The system operates strictly locally, eliminating the need for complex cloud orchestration components and decreasing latency natively.

\begin{figure}[htp]
    \centering
    % Placeholder for the architecture diagram which we will generate as a flowchart
    \begin{tikzpicture}[node distance=2cm and 2cm,
    box/.style={rectangle, draw, minimum width=3cm, minimum height=1cm, text centered, rounded corners},
    arrow/.style={thick,->,>=stealth}]

    \node (user) [box, fill=blue!10] {Administrator};
    \node (engine) [box, fill=red!10, below left=of user] {C++ Execution Engine};
    \node (ui) [box, fill=green!10, below right=of user] {Python Admin Panel};
    
    \node (policy) [box, fill=yellow!10, below=of engine] {Policy Enforcer};
    \node (session) [box, fill=yellow!10, right=of policy] {Session Manager};
    
    \node (iptables) [box, fill=gray!10, below=of policy] {IP Tables / FS Mounts};
    \node (os) [box, fill=gray!10, below=of session] {Linux OS Commands};

    \draw [arrow] (user) -- node[anchor=east, align=center] {vanish start/stop} (engine);
    \draw [arrow] (user) -- node[anchor=west, align=center] {Views UI \\ \& Presets} (ui);
    \draw [arrow, dashed] (ui) -- node[anchor=north] {Reads Logs/Configs} (engine);
    
    \draw [arrow] (engine) -- (policy);
    \draw [arrow] (engine) -- (session);
    
    \draw [arrow] (policy) -- (iptables);
    \draw [arrow] (session) -- (os);
\end{tikzpicture}

    \caption{Vanish Lite High-Level Architecture Diagram}
    \label{fig:architecture}
\end{figure}

\section{Core Components}
\subsection{C++ Session Engine}
The engine accepts commands via a standard CLI interface (\texttt{main.cpp}). It parses the requested operational mode and directly manipulates the system's underlying \texttt{passwd} user databases to generate sandboxed temporary users (\texttt{vanish\_*}). State management is orchestrated through the \texttt{session\_manager.cpp}, which stores live session timers within protected directories (\texttt{/var/vanish\_sessions/}). 

\subsection{Policy Enforcer}
Located fundamentally in \texttt{policy\_enforcer.cpp}, this component dynamically generates Linux network rules applying IP tables blocks and routing overrides (such as overriding DNS files per user). Depending on the \texttt{Strict Mode} or \texttt{Online Mode}, specific URL allow/block combinations are enforced immediately upon session instantiation.

\subsection{Python Admin Dashboard}
Built utilizing standard Python synchronous methodologies, the API parses static log structures and config snapshots to render a visual representation of active sessions, security overrides, and system compliance. The dashboard includes a dry-run feature natively simulating the engine parameters before executing root-level modifications.
